\section{Background}
\label{sec:background}
\paragraph{Notifications and durable state.} \texttt{ContentObserver} callbacks require a live registered process~\cite{android_contentobserver}. ContentService immediately serves sufficiently important processes or writes marked \texttt{NOTIFY\_NO\_DELAY}; other callbacks wait its 10-second background delay~\cite{aosp_contentservice}. JobScheduler content-URI triggers survive process death but cannot persist across reboot~\cite{aosp_jobinfo}. Neither mechanism provides a replayable provider history. Counters, versions, timestamps and tombstones may supply that history if their coverage and reset contracts hold; direct reads instead obtain current query results.

\paragraph{Provider and lifecycle boundaries.} Contacts and Calendar sync metadata belongs to the account's sync adapter, which clears dirty/deleted flags after synchronization~\cite{android_syncadapter}. A third-party reader cannot assume those flags constitute its own change log. Android exposes MediaStore generation evidence, while iOS Contacts and Photos expose reader-side history~\cite{apple_cnchangehistory,apple_phpersistent}.

Cached processes may be killed or frozen~\cite{lebeck2020marvin,liang2020acclaim,android_freezer}. One-way observer calls wait during freezing. Force-stop cancels jobs and alarms; permission revocation kills the process; unregistration removes callback coverage. Doze and standby buckets constrain background execution~\cite{android_doze}; our tested schedules show no additional observer delay beyond background deferral. A new process identity or \texttt{ApplicationExitInfo} identifies absence, not intervening provider changes~\cite{android_exitinfo}.

\paragraph{Scope.} We study benign staleness with recorded sources and queries. Lifecycle probes use deterministic scripts; derived-fact and resumed-task experiments invoke local LLMs. The target is external-state memory, not planning quality. Screen-only facts without a supported query binding remain \Unknown.
